% Pursuing Stacks (cotes 134-2 to 134-8): archive cote 134-3, site page 42 (PDF page 43), typescript page 42.
% typescript. Draft compared: ps2.tex:2058-2133.
% AI-transcribed by Claude Opus 5.5 (claude-opus-5-5, adaptive thinking, effort high), Message Batches API, batch ps1 (msgbatch_01Rv5hCaNTtUt4JHx1yqoHPL), 2026-11-12; not yet checked by eye.
\scanpage{134-3/42}
\mspage{42}

15.3.

32.   I definitely have the feeling to be out of the thin air  -  the con-%
ditions  (T 1) to (T 3) look to me so elegant and convincing,  that I have
no doubts left they are ``the right ones'' !  A lot of things comes to mind
what to do next,  I'll have to look at them one by one though. Let me make
a quick provisional planning:

1)  Have a closer look at the conditions  T 1) to  T 3),  to see how fa
% NOTE: the right edge of the scan cuts off the end of this line and of the next
they are necessary for \xout{the} $A$ to be a test category in the (admittedly provi-%
sional)sense\xout{\unsure{;}} I \xout{used this word}\inserted{gave to this notion} last week and yesterday,  and to pin down
the feeling of these being just the right ones.

2)  Use these conditions for constructing lots of test categories,
including all the simplicial and cubical types which have been used so far,

3)  Check that these conditions are stable under taking the product of
two or more test categories,  and possibly use this fact for comparing
the homotopy theories defined by any two such categories.

4)  Look up \xout{if} (using T 1 to T 3) if an elementary \deleted{\unsure{topos}} modelizer $A^{\wedge}$
is indeed a ``closed model category'' in Quillen's sense,  and maybe too get
a better feeling of how far apart (Cat) is from being a closed model cate-%
gory\hand{(9)}. Visibly there \uline{are} some natural constructions in homotopy theory whi\unsure{\hand{ch}}
do make sense in (Cat).

5)  Using the understanding of test categories obtained,  come back to
the question of which categories M associated to \xout{an} algebraic structure
species can be viewed as modelizers,  and to the question of unicity or ca-%
nonicity of the modelizing structure $W\hand{$\subset$}F\over{\unsure{1}}{l}(M)$.

T\over{\unsure{n}}{h}at makes a lot of questions to look at,  and the theory of stacks I
set out to sketch seems to be fading ever more into the background\over{.}{!}  It is
likely though that a better \inserted{general} understanding of the manifold \inserted{ways of} constructi\over{ons}{ng} \deleted{\unsure{of}}
the category (Hot) of homotopy types will not be quite useless,  when get-%
ting back to the initial program,  namely stacks.  Quite possibly too,  on
my way I will have to remind myself \inserted{of} and look up,  in the present setting,
\deleted{of} the main structural properties of (Hot),  including exactness properties\hand{,}
and generators and cogenerators.  This also reminds me of an \xout{old} intriguing
foundational question since the introduction of derived categories and their
non commutative analogs,  which I believe has never been settled yet,  namel%
\hand{-}ly the following:

6)  In an attempt to grasp the main natural structures associated to
derived categories,  and Quillen's \xout{n}non commutative analogons including
(Hot),  try to develop  a comprehensive notion of a ``triangulated category''\hand{,}
without the known drawbacks of \xout{the notion} Verdier's provisional notion.
